\section{Introduction}
\label{sec:introduction}

Robotic wiping is a contact-rich control problem in which geometric progress
and force regulation must succeed together. The tool must first acquire the
surface, approach a requested load without a large transient, maintain useful
contact during motion, and recover when contact is lost. These phases place
different demands on a controller. Contact acquisition rewards decisive
inward motion, stable tracking rewards accurate force regulation, and rapid
force growth calls for conservative action. A planning objective and sampling
budget that remain fixed throughout an episode need not serve all three
requirements equally well.

Temporal-difference model-predictive control (TD-MPC2) offers a natural basis
for this problem: a learned actor provides a policy prior, while a latent
world model evaluates action sequences online
\cite{hansen2024tdmpc2,williams2018information}. Earlier ForceWipe experiments
showed that behaviour-cloning (BC)-assisted TD-MPC2 could complete controlled
first-pass wiping directly at 5, 8, and 12~N. They also exposed three limits.
First, a hard measured-force threshold routed between the actor and planner
without adapting the planning objective itself. Second, the 12~N tier retained
a systematic negative force bias. Third, the matched actor-only comparison
did not resolve a general performance benefit from online planning despite a
clear computation cost. These observations motivate a method that changes
what the planner optimizes and how much it searches, rather than merely
changing when planning is enabled.

We introduce regime-adaptive risk-aware MPPI (RA-RMPPI), a unified planner for
direct TD-MPC2 wiping. The requested force is encoded by a dedicated condition
branch and retained through imagined latent transitions. Learned next-force
predictions provide a causal risk signal. Smooth causal memberships represent
acquisition, approach, tracking, recovery, and transient risk. They
continuously adjust contact, tracking, transient, and actor-anchor costs
together with the MPPI sample count, optimization
iterations, and exploration scale. The actor remains the proposal
distribution in every regime; the historical hard actor--planner switch is
removed.

The empirical design separates method selection from evaluation. Five
unfiltered seeds share the same data and update budget. A four-arm ablation
isolates adaptive objectives and adaptive computation, while additional
ablations examine behaviour-cloning strength, transient-risk representation,
and the historical routing threshold. The method is then evaluated on new
factor-separated shifts in surface geometry and contact dynamics, avoiding
the path--surface--tool confounding of the earlier combined
out-of-distribution (OOD) blocks. Across 270 paired evaluations, the complete
method raises compound pass from 85/135 to
100/135 while using fewer planning samples and less host planning time. The
gain is concentrated in force regulation: task completion and sampled force
violations do not improve, and the cylindrical 12~N condition remains a
failure boundary. A zero-shot MuJoCo stress test sharpens that boundary: all
30 flat and inclined evaluations achieve the compound outcome, whereas none of
the 15 cylindrical evaluations does.

The paper makes three contributions. First, it presents a force-conditioned
latent dynamics construction that preserves the requested force during
multi-step imagination. Second, it develops a regime-adaptive risk objective
and compute schedule that unifies acquisition, tracking, and recovery without
a discrete actor--planner router. Third, it provides matched component and
training ablations, including negative evidence for a more elaborate
transient-envelope representation, together with factor-separated and
cross-engine simulation evaluation. The resulting evidence identifies a
statistically resolved tracking--computation improvement while preserving the
method's unresolved curvature and force-tail limits.
